% SOURCE: original Chinese PDF, viewer pages 7--9.
\section{Eigenvalues and Eigenvectors of Operators (Matrices)}
\subsection*{1) Eigenvalues and eigenstates of a general operator}
The eigenvalue equation of an operator is
\begin{equation}
\op T\ket a=\lambda\ket a,
\end{equation}
where $\lambda$ is called an eigenvalue of $\op T$ and $\ket a$ is called an eigenvector of $\op T$.

In matrix form (prove this yourself using the completeness relation),
\begin{equation}
Ta=\lambda a,\qquad (T-\lambda I)a=0.
\end{equation}
The condition for a nonzero eigenvector $a\ne0$ is
\begin{equation}
\det(T-\lambda I)=0.
\end{equation}
This is the secular equation in an $N$-dimensional space:
\begin{equation}
\begin{vmatrix}
T_{11}-\lambda&T_{12}&\cdots&T_{1N}\\
T_{21}&T_{22}-\lambda&\cdots&T_{2N}\\
\vdots&\vdots&\ddots&\vdots\\
T_{N1}&T_{N2}&\cdots&T_{NN}-\lambda
\end{vmatrix}=0.
\end{equation}
One thereby obtains $N$ eigenvalues $\lambda_i$ ($i=1,2,\ldots,N$). Substituting each into
\begin{equation}
(T-\lambda_i I)a=0
\end{equation}
gives the corresponding eigenvector $a$.

\textbf{Example.} Find the eigenvalues and eigenvectors of
\begin{equation}
T=\begin{pmatrix}1&0\\0&-1\end{pmatrix}.
\end{equation}
The secular equation is
\begin{equation}
\begin{vmatrix}1-\lambda&0\\0&-1-\lambda\end{vmatrix}=0,
\quad\text{i.e.,}\quad \lambda^2-1=0,
\end{equation}
and the eigenvalues are $\lambda_\pm=\pm1$.

Let $a=(a_1,a_2)^{\mathsf T}$. For $\lambda=1$,
\begin{equation}
\begin{pmatrix}0&0\\0&-2\end{pmatrix}
\begin{pmatrix}a_1\\a_2\end{pmatrix}=0.
\end{equation}
The normalized eigenvector is $a=(1,0)^{\mathsf T}$. Similarly, taking $\lambda=-1$ gives $a=(0,1)^{\mathsf T}$.
\subsection*{2) Hermitian operators}
If $\op T^\dagger=\op T$, or $T^\dagger=T$, the operator $\op T$ is called Hermitian and $T$ is called a Hermitian matrix.

We discuss the following properties of Hermitian operators.
\paragraph{a) Conjugate symmetry.}
\begin{equation}
\mel a{\op T}b=\mel b{\op T^\dagger}a^*=\mel b{\op T}a^*.
\end{equation}
In particular,
\begin{equation}
\mel a{\op T}a=\mel a{\op T}a^*,
\end{equation}
which shows that the expectation value $\mel a{\op T}a$ of a Hermitian operator is real.

For an anti-Hermitian operator, $\op T^\dagger=-\op T$ and
\begin{equation}
\mel a{\op T}a=-\mel a{\op T}a^*,
\end{equation}
so its expectation value is purely imaginary.
\paragraph{b) Reality of eigenvalues and orthogonality.}
Suppose
\begin{equation}
\op T\ket i=\lambda_i\ket i,\qquad \op T\ket j=\lambda_j\ket j.
\end{equation}
From
\begin{equation}
\mel j{\op T}i=\lambda_i\braket j i
\end{equation}
and
\begin{equation}
\bra j\op T=\bra j\lambda_j^*,\qquad
\mel j{\op T}i=\lambda_j^*\braket j i,
\end{equation}
we obtain
\begin{equation}
(\lambda_i-\lambda_j^*)\braket j i=0.
\end{equation}
For $j=i$ and $\braket i i\ne0$, $\lambda_i=\lambda_i^*$: the eigenvalues of a Hermitian operator are real.

For $j\ne i$, if $\lambda_i-\lambda_j\ne0$, then $\braket j i=0$. Thus eigenstates corresponding to different eigenvalues of a Hermitian operator are orthogonal. Since they can always be normalized, the orthonormality condition is
\begin{equation}
\braket i j=\delta_{ij}.
\end{equation}
\textbf{Question:} Are eigenstates corresponding to the same eigenvalue orthogonal?
\paragraph{c) Orthogonalization by linear superposition (Schmidt orthogonalization).}
Suppose several eigenstates correspond to the same eigenvalue, i.e., there is degeneracy. For example, for a $g$-fold degenerate eigenvalue ($g\ge2$),
\begin{equation}
\op T\ket{i,j}=\lambda_i\ket{i,j},\qquad j=1,\ldots,g.
\end{equation}
Are these $g$ eigenvectors orthogonal?

Redefine $g$ new states:
\begin{equation}
\ket{i,n}=\sum_{j=1}^g C_{nj}\ket{i,j},\qquad n=1,2,\ldots,g.
\end{equation}
Since
\begin{equation}
\op T\ket{i,n}=\sum_{j=1}^g C_{nj}\op T\ket{i,j}
=\lambda_i\sum_{j=1}^g C_{nj}\ket{i,j}=\lambda_i\ket{i,n},
\end{equation}
each $\ket{i,n}$ is still an eigenstate of $\op T$ with eigenvalue $\lambda_i$. Hence it suffices to show that the $\ket{i,n}$ can be made orthonormal.

Can one choose the coefficients $C_{nj}$ so that these $g$ new states satisfy
\begin{equation}
\braket{i,m}{i,n}=\delta_{mn}\;?
\end{equation}
There are $g$ normalization equations and $(g^2-g)/2$ orthogonality equations, making $g(g+1)/2$ independent equations, fewer than the $g^2$ coefficients $C_{nj}$. Thus there are various choices of coefficients satisfying the orthonormality conditions, so that the new states $\ket{i,n}$ are orthonormal:
\begin{equation}
\braket{i,m}{j,n}=\delta_{ij}\delta_{mn}.
\end{equation}
Here $\delta_{ij}$ follows from orthonormality of eigenstates belonging to different eigenvalues, while $\delta_{mn}$ follows from orthogonalization by linear superposition.

\textbf{Conclusion:} Eigenstates of a Hermitian operator can be chosen orthonormal, whether the eigenvalues are degenerate or nondegenerate.
\paragraph{d) Completeness.}
It can be shown that eigenvectors of a Hermitian operator satisfy
\begin{equation}
\sum_i\ket i\bra i=\mathbf1.
\end{equation}
Hence the eigenvectors of a Hermitian operator can form a set of orthonormal basis vectors of a Hilbert space, i.e., a linear vector space or a representation.